\begin{claim}{D.37}{Proved}{Negative transpositions and the specified four-cusp deck action}
Every involution of the full three-cusp loop-outer group whose
puncture permutation is a transposition and whose multiplier has
$\lambda_2=-1$ gives a faithful $V_4$ action on a specified
four-cusp double cover.  One generator is the standard free positive
deck involution; another has double-transposition permutation and
multiplier $\lambda$, and its center-free extension is torsion-free.
Conversely every such $V_4$ action containing the specified standard
deck action descends to a negative transposition on the three-cusp
group.  The constructions are inverse as marked extensions.

For the double-transposition permutation of D.29, neither of the
commuting single transpositions preserves any separating two-pants
partition whose pants the double transposition exchanges.  Thus an
additional deck action in this construction necessarily moves the
separator.  This is an obstruction to preserving that decomposition,
and asserts no obstruction to an action moving it.
\end{claim}
\begin{proof}
We specify both the markings and all the extension maps.  Put
\[
 F=\widehat F(x,y),\qquad z=(xy)^{-1},\qquad
 \chi:F\longrightarrow C_2,\quad (\chi(x),\chi(y))=(0,1),
 \qquad N=\ker\chi.
                                                        \tag{D.37a}
\]
Schreier rewriting identifies $N$ with the full free profinite
group on
\[
 A_0=x,\qquad A_1=yxy^{-1},\qquad B=y^2,
 \qquad D=z^2=B^{-1}A_1^{-1}A_0^{-1}.
                                                        \tag{D.37b}
\]
Its four peripheral generators satisfy $A_0A_1BD=1$.
These are the actual degree-one lifts of $x$ and the degree-two
lifts of $y,z$; completion of the discrete index-two Schreier
isomorphism gives precisely the specified inclusion $N\subset F$.
Let $\delta\in\Out(N)$ be conjugation by $y$.  Its named
representative is
\[
 A_0\longmapsto A_1,\qquad
 A_1\longmapsto B A_0 B^{-1},\qquad B\longmapsto B,
 \qquad \delta^2=\operatorname{Inn}(B).
                                                        \tag{D.37c}
\]
Here the final equality concerns the named automorphism;
its outer class has order two, multiplier $1$, and permutation
$(A_0\ A_1)$.  Its center-free extension is the explicitly
identified group $F$, hence is torsion-free.  In the converse
below, ``the specified deck action'' means this action transported
by a given peripheral isomorphism.  It does not mean an arbitrary
positive outer involution with a torsion-free extension.

Let $g$ be a negative transposition, labeled to fix $x$ and
exchange $y,z$.  Its center-free extension $E$ contains $F$
as an open subgroup of index two.  The peripheral normalizer
surjects onto $\langle g\rangle$, and A.7 gives its kernel
$\overline{\langle x\rangle}=\widehat\Z$.
For a lift $v$ in this normalizer write
\[
 v x v^{-1}=x^\lambda,\qquad v^2=x^a,
 \qquad (\lambda-1)a=0,\qquad \lambda^2=1.
                                                        \tag{D.37d}
\]
At every prime where $\lambda_p=-1$ the last commutation
identity gives $a_p=0$.  At every prime where $\lambda_p=1$
the prime is odd, since $\lambda_2=-1$, so $2$ is invertible.
Choose $b_p=-a_p/2$ in these latter coordinates and $b_p=0$
in the former.  Then $s=x^b v$ has order two.  Thus
\[
 E=F\rtimes_\alpha\langle s\rangle,\qquad
 \alpha^2=1,\qquad \alpha(x)=x^\lambda.
                                                        \tag{D.37e}
\]
This constructs the required finite-order lift inside the
peripheral normalizer, without a further choice of inner
normalization.  The multiplier and the cusp permutation give
$\chi\alpha=\chi$.  Therefore $N$ is normal in $E$, and
$E/N=C_2\times C_2$, with basis the images of $y,s$.

There is exactly one $F$-conjugacy class of order-two subgroups
of $E$.  Indeed the action on $H^1(F,\F_2)$ is the two-dimensional
nontrivial unipotent module $J_2$, isomorphic to the regular
$\F_2[C_2]$ module.  Consequently
$H^i(C_2,H^1(F,\F_2))=0$ for $i>0$.
The Hochschild--Serre spectral sequence, with the split
inflation on its bottom row, gives
\[
                 \dim H^n(E,\F_2)=1\qquad(n\geq3).
                                                        \tag{D.37f}
\]
The full profinite Brown restriction theorem used in D.28
identifies this high-degree cohomology with the product over
the order-two conjugacy classes.  Each class contributes
$1+\dim H^1(C_F(T),\F_2)$, since its normalizer is
$C_F(T)\times C_2$; see
\cite[Theorem 2 and \S1.1]{A-KochloukovaZalesskii2007}.
The section in (D.37e) exists, so (D.37f) forces exactly one
class and zero first mod-two cohomology for its centralizer.
Conjugacy by $E$ agrees here with conjugacy by $F$: for any
section $T$, one has $E=FT$, and the $T$ factor centralizes
that subgroup itself.
Every finite-order element of $E$ injects into $E/F$,
because $F$ is torsion-free, and consequently has order two.
All nonidentity torsion therefore maps to the one line
$\langle\bar s\rangle$ in the abelian group $E/N$.

Conjugation gives a faithful action of $E/N$ on $N$ in
$\Out(N)$.  To check the kernel, first
$C_F(N)=1$: an element centralizing both $x$ and $y^2$
lies in $\langle x\rangle\cap\langle y\rangle=1$ by A.7.
For $e\in C_E(N)$ and $f\in F$, the commutator $[e,f]$
lies in $F\cap C_E(N)$ and hence is trivial.  Thus $e$
centralizes $F$; faithfulness of the original nonidentity
outer action and $Z(F)=1$ give $e=1$.

Write $a$ for the outer restriction of $\alpha$ to $N$.
It fixes the classes of $A_0,A_1$ and exchanges $B,D$.
For the second fixed class use the explicit conjugator
$\alpha(y)y^{-1}\in N$:
\[
 \alpha(A_1)
  =(\alpha(y)y^{-1}) A_1^\lambda
                         (\alpha(y)y^{-1})^{-1}.
                                                        \tag{D.37g}
\]
The unique lifts of the other two cusps are exchanged inside
$N$.  Thus $\delta$ and $a$ commute in $\Out(N)$, and
\[
 \tau=\delta a,\qquad
 \sigma(\tau)=(A_0\ A_1)(B\ D),\qquad
 \operatorname{or}(\tau)=\lambda.
                                                        \tag{D.37h}
\]
The preimage $E_\tau$ of the diagonal line
$\langle\bar y+\bar s\rangle$ in $E/N$ is the center-free
extension of this outer involution.  It is torsion-free:
all nonidentity torsion has image $\bar s$, and its
intersection with $N$ is already torsion-free.
For an explicit representative set
\[
 \beta=\operatorname{Inn}(y)\alpha|_N,\qquad
 c=y\alpha(y)\in N;\qquad
 \beta^2=\operatorname{Inn}(c),\quad \beta(c)=c.
                                                        \tag{D.37i}
\]
Then $E_\tau$ is exactly the twisted extension generated by
$N$ and $ys$, whose square is $c$; the multiplication rule
is the rule of (D.29.9) with these named $\beta,c$.

Conversely let a marked four-cusp group carry the specified
$\delta$ and an involution $\tau$ with the permutation and
multiplier in (D.37h), satisfying $[\delta,\tau]=1$.
The distinct permutations make
$\Gamma=\langle\delta,\tau\rangle\cong V_4$.
Take its center-free extension
\[
 K=\{(\varphi,\gamma)\in\Aut(N)\times\Gamma:
                         [\varphi]=\gamma\}.
                                                        \tag{D.37j}
\]
Here $N$ embeds by its inner automorphisms.  The preimage
$K_\delta$ of $\langle\delta\rangle$ is identified with
the exact group $F$ in (D.37a) by
\[
 f\longmapsto
       (\operatorname{Inn}(f)|_N,\delta^{\chi(f)}).
                                                        \tag{D.37k}
\]
Injectivity follows from $C_F(N)=1$; surjectivity follows
because an automorphism in either of these two outer
classes differs from $1$ or $\operatorname{Inn}(y)|_N$
by an inner automorphism from $N$.
The quotient $K/F=C_2$ therefore gives an outer action on
this very $F$.  It is faithful, since any element of $K$
centralizing $F$ centralizes its inner copy of $N$, and
an automorphism of the center-free $N$ with that property
is the identity.

We check descent of the peripheral powers, not only of the
four unordered procyclic subgroups.  The two classes
$A_0,A_1$ become the one $F$-class of $x$.
The other classes are $B=y^2$ and $D=z^2$.
For a representative of $\tau$, the induced automorphism
$\psi$ of $F$ satisfies
$\psi(y)^2=h z^{2\lambda}h^{-1}$ for some $h\in N$.
A.7 implies
$C_F(hz^{2\lambda}h^{-1})=h\langle z\rangle h^{-1}$,
so $\psi(y)=hz^\mu h^{-1}$ and $2\mu=2\lambda$.
Multiplication by two is injective on $\widehat\Z$;
hence $\mu=\lambda$.  The same proof applies to $z$,
and the $x$ exponent is already given by $A_0$.
Thus the descended action fixes $x$, exchanges $y,z$,
and has exactly the original $\lambda$, in particular
$\lambda_2=-1$.

For completeness its normalizer construction (D.37d)--(D.37e)
and (D.37f) also recover the torsion-free assertion upstairs,
without making that assertion an extra converse hypothesis.
An order-two lift in the normalizer of $\langle x\rangle$
fixes $A_0$, so its image in $\Gamma$ is $\delta\tau$,
not $\tau$.  All torsion has this same image by (D.37f).
The preimage of $\langle\tau\rangle$ is therefore torsion-free.
The conjugation maps (D.37j)--(D.37k) identify the recovered
extension with the original one on every element; conversely
restriction to $N$ recovers the given $\Gamma$ and both
named generators.  This proves the asserted inverse
correspondence of marked extensions.

We finish by making the obstacle for the D.29 gluing explicit.
Write its ordered cusps as $1,2,3,4$, so that
$\sigma(\tau)=(1\ 3)(2\ 4)$.  The unordered two-pants
partitions and their images are
\[
\begin{array}{c|c|c}
 \text{partition}&\tau\text{ on its two blocks}
                     &(1\ 3)\text{ or }(2\ 4)\\\hline
 12\mid34&\text{exchanges}&14\mid23\\
 13\mid24&\text{preserves each}&13\mid24\\
 14\mid23&\text{exchanges}&12\mid34.
\end{array}                                            \tag{D.37l}
\]
The only single transpositions commuting with
$(1\ 3)(2\ 4)$ are $(1\ 3)$ and $(2\ 4)$:
a commuting permutation must send its two two-cycles to
two-cycles, and a transposition can do this only by
exchanging the members of one of them.  Hence a deck
involution of the required permutation cannot act on the
same two-vertex pants graph as the inverse-pair gluing,
under any relabeling or choice of its exchanged partition.

One can nevertheless write the required geometric deck
map in the D.29 free basis.  For
$(u_1,u_2,u_3,u_4)=(a,a^{-1}l,b,b^{-1}l^{-1})$, the
change of four-cusp marking
\[
 A_0=a,\quad A_1=(a^{-1}l)b(a^{-1}l)^{-1},\quad
 B=a^{-1}l,\quad D=b^{-1}l^{-1}
                                                        \tag{D.37m}
\]
satisfies $A_0A_1BD=1$.  Its inverse substitutions are
$a=A_0$, $l=A_0B$, $b=B^{-1}A_1B$.
Transporting (D.37c) through these maps gives
\[
 \Delta(l)=a^{-1}lb,\qquad
 \Delta(a)=a^{-1}lbl^{-1}a,\qquad
 \Delta(b)=a,\qquad
 \Delta^2=\operatorname{Inn}(a^{-1}l).
                                                        \tag{D.37n}
\]
It has multiplier $1$ and permutation $(1\ 3)$.
The existence question for D.29 is the existence of a
deck action of this marked geometric kind commuting
outerly with its actual $\alpha$; (D.37l) alone does
not decide it.

In particular failure of commutation for one selected
deck representative cannot settle the question.  In the
geometric specialization $\lambda=-1$, $\epsilon=1$,
$\kappa=1$, $f=1$, D.29 gives
\[
 \alpha(l)=l,\quad\alpha(a)=b^{-1},\quad
                         \alpha(b)=la^{-1}l^{-1}.
\]
For (D.37n), the words
$\alpha\Delta(l)=bl^2a^{-1}l^{-1}$ and
$\Delta\alpha(l)=a^{-1}lb$ are not conjugate even
in the full profinite group: the finite quotient
$l\mapsto(1\ 2\ 3)$, $a,b\mapsto(2\ 3)$ in $S_3$
sends the first to $1$ and the second to a nonidentity
three-cycle.  These two automorphisms therefore do not
commute outerly.  Yet another geometric deck does commute:
\[
 \begin{aligned}
 \Delta_0(l)&=bla^{-1},\qquad \Delta_0(a)=b,\qquad
                  \Delta_0(b)=blal^{-1}b^{-1},\\
 \Delta_0^2&=\operatorname{Inn}(bl),\qquad
       \alpha\Delta_0=\operatorname{Inn}(b^{-1})\Delta_0\alpha.
 \end{aligned}
                                                        \tag{D.37o}
\]
Its origin is also explicit.  In the geometric extension
$E=\widehat F(a,t)$, put $l=t^2$ and
$b=ta^{-1}t^{-1}$.  The involution
$\eta(a)=a^{-1}$, $\eta(t)=a^{-1}t$ preserves the
index-two kernel.  Restricting
$\operatorname{Inn}(t)\eta$ to that kernel gives
exactly (D.37o).  The formulas prove its peripheral
permutation and multiplier as well as the two displayed
composition identities.  Thus this geometric
specialization supplies an explicit example of the
separator being moved.
To verify directly that $\Delta_0$ is a deck of the
specified geometric kind, adjoin its lift $r$ with
$r^2=bl$.  The extension presentation eliminates to
$\widehat F(a,r)$ by
$b=rar^{-1}$ and $l=b^{-1}r^2=ra^{-1}r$.
The peripheral generators downstairs are
$a,a^{-1}r,r^{-1}$, whose product is one.
Their lifts are the two classes of $a,b$, the square
$(a^{-1}r)^2=a^{-1}l$, and the square $r^{-2}$,
conjugate to $b^{-1}l^{-1}$.  These substitutions identify
its four-cusp inclusion with (D.37a)--(D.37b).

The executable \texttt{code/part\_c/negative\_transposition\_d37.py}
verifies the free-word covering and inverse substitutions,
the deck squares, all the equations in (D.37o), the
two-dimensional mod-two cohomology calculation, and the
complete permutation calculation (D.37l) under all
relabelings.  The extension and torsion arguments above
prove the statements for the full profinite groups.
\end{proof}
